We report Monte Carlo experiments on rank-$r$ LoRA-like
ensembles validating the theoretical rates and constants. All
runs use the hard distribution $P^\star$ of \S\ref{sec:setup}
with task-dependent spectra $(D_t)_{t=1}^T$. Code, random seeds,
and raw outputs are released with the paper.

\subsection{Synthetic validation}\label{subsec:exp-synthetic}\label{subsec:exp-T-synthetic}

\paragraph{Lower-bound slope.} We sample ensembles across
$T\in\{2,3,4,8\}$, $r\in\{4,8,16\}$, $d\in\{128,512,2048\}$,
with both isotropic spectra ($D_t = I_r$) and four
heavy-anisotropy families (geometric, linear,
two-scale, and mixed; spectra spread across
$[2^{-1}, 2^1]$). For each configuration we measure the excess
$\mathbb{E}_{P^\star}[\|\bar H^{1/2}w^\star - \eta\|^2]$ of the
Gaussian-QR + uniform scalar quantization algorithm of
\S\ref{sec:achv} as a function of $R = b\deff$ bits, and fit a
linear regression of $\log_2\mathbb{E}[\text{excess}]$ against
$b = R/\deff$. Across all tested configurations the empirical
slope is $-2.00 \pm 0.01$ (95\% bootstrap CI across 1000 trials
per configuration), matching the prediction $-2R/\deff$ of
Theorem~\ref{thm:general}. The floor term $B^2(1 - \E[\deff]/
(Tr))$ is reproduced to four to five digits of precision in
the common-$D_0$ and task-dependent-$D_t$ regimes.

\paragraph{Achievability constant and its $T$-dependence.} We measure
the ratio
$\mathbb{E}_{P^\star}[\max_t D_t(w^\star)] / (c_{\mathrm{TQ}}
B^2 2^{-2R/(Tr)})$ of empirical upper bound to the Stiefel
lower bound of Theorem~\ref{thm:general} at $c_{\mathrm{TQ}} = 1$.
The ratio is constant in $R$ — a necessary condition for the UB
to match the LB in order — and starts near $11$ at $T = 2$
($r = 4$, $d = 128$, isotropic and task-mismatched anisotropic
spectra), consistent with $C = Tc^2/3 \approx 17$ at $c = 5$ combined
with $c_{\mathrm{TQ}} \approx 1$ from
TurboQuant~\cite{zandieh2025turboquant}. Sweeping
$T \in \{2, 4, 8, 16\}$ ($1000$ trials/cell) settles the
$T$-dependence flagged in Remark~\ref{rem:tgap}: the constant grows
\emph{logarithmically}, not linearly --- a $\log T$ fit gives
$R^2 \geq 0.998$ against $R^2 \approx 0.91$ for a linear fit, and the
$T{=}16$ value is only $1.77\times$ the $T{=}2$ value, far below the
$8\times$ a strictly linear-$T$ factor would give. The linear $T$ in
$C = Tc^2/3$ is thus an artifact of the crude
$\max_t \leq T\,\mathrm{avg}$ step (Remark~\ref{rem:tgap}), not
intrinsic to the rate-distortion function; the matching $\log T$
achievability constant is stated as a conjecture.

\paragraph{Shared-$V$ regime.} For the floor-positive regime
we implement the null-space-aware bit allocation of
\S\ref{subsec:sharedv} with per-coord fractional bits and a
locked clip $c = 11.5\sigma_{\mathrm{pc}}$ (tuned once, applied
across all configurations). At $T = 2$, $r = 4$, across five
anisotropy mixes (iso+iso, iso+geom, geom+twoscale, lin+twoscale,
geom+iso), the empirical slope of excess over
$\mathrm{cheb}^2$ is $-1.60 \pm 0.10$ at $n_{\mathrm{trials}} =
1000$, matching the theoretical prediction
$-2r/(r + |A| - 1) = -8/5$ with bootstrap 95\% CI $\pm 0.010$
per cell. The $\pm 0.10$ spread reflects structural anisotropy:
iso+iso tracks closer to the scalar-quantization limit $-2$;
geom+twoscale undershoots due to $H_t$-metric distortion on the
perpendicular subspace. Extending to $T = 3, 4$ via a SOCP
Chebyshev solver with Gauss-Newton KKT refinement (residual
$\sim 10^{-15}$), the null-space split continues to track the
theoretical prediction across mixed anisotropies.

\paragraph{Random-codebook comparison.} As a sanity check on
the tightness of our linear construction in the shared-$V$
regime, we ran a valid rate-$R$ random-codebook encoder ($2^R$
iid Gaussian codewords, encoder picks the one minimizing
$\max_t D_t$, $R \leq 18$). At $T = 3$, $r = 3$, the
random-codebook slope is $-1.45$, strictly steeper than our
explicit linear construction's $-1.20$. This rules out
sharpening the lower bound to match the linear upper-bound
exponent; the true rate-distortion function in the
floor-positive regime lies strictly between $-2r/r = -2$ and
$-2r/(r + |A| - 1)$ and is left open
(Remark~\ref{rem:sharedv-gap}).

\subsection{Real-LLM merging}\label{subsec:exp-lora}\label{subsec:exp-multiseed}

We calibrate the bound against real LoRA fine-tunes across four
instruction-tuned base models and four tasks. The aim is not to
reproduce the idealized slopes of \S\ref{subsec:exp-synthetic} --
real task vectors are not drawn from the worst-case distribution
$P^\star$ -- but to test which \emph{qualitative} predictions of
the theory survive contact with real models, and to quantify how
far current merging methods operate from the
information-theoretic floor.

\paragraph{Setup.} We use four base models spanning distinct
pretraining recipes and two attention designs:
Llama-3.1-8B-Instruct (multi-head attention, MHA),
Qwen-2.5-7B-Instruct (grouped-query attention, GQA),
Mistral-7B-Instruct-v0.3 (MHA), and Yi-1.5-9B-Chat (MHA,
Llama-family architecture). For each model we train $T=4$
task-specific LoRA adapters (rank $r = 16$ on the
$\{q,k,v,o\}$ projections of every attention block; AdamW,
learning rate $5\times10^{-5}$, $7500$ examples $\times$ $2$
epochs, full bf16 base) on GSM8K (math), Alpaca-cleaned
(instruction following), Magicoder-OSS-Instruct (code), and
WMT19 en$\to$de (translation). This yields $16$ adapters in
total. We evaluate on $n = 1000$ held-out examples per task,
disjoint from training, and report negative log-likelihood (NLL)
in nats per response token. Confidence intervals are $95\%$
percentile bootstrap over $10{,}000$ resamples of the per-example
NLL.

\paragraph{Distortion metric.} For a merged delta $w$ we measure
the per-task \emph{excess} NLL over the task's own adapter
$\tau_t$, $\Delta L_t(w) = L_t(\theta_0 + w) - L_t(\theta_0 +
\tau_t)$, and the worst-task distortion $\widehat D(w) = \max_t
\Delta L_t(w)$, the operational analogue of the
$\max_t$ objective in \S\ref{sec:setup}. The shared frozen base
$\theta_0$ cancels in $\Delta L_t$, so the metric is invariant
to base-model quality and isolates the cost of merging.
Table~\ref{tab:nll-tau} reports the per-adapter baselines
$L_t(\theta_0 + \tau_t)$.

\begin{table}[t]
\centering
\caption{Per-adapter baseline NLL $L_t(\theta_0+\tau_t)$ (nats/token,
$n=1000$). Lower is better. Yi-1.5-9B is the strongest per-task
learner on three of four tasks.}
\label{tab:nll-tau}
\begin{tabular}{lrrrr}
\hline
Model & GSM8K & Alpaca & Magicoder & Translation \\
\hline
Llama-3.1-8B    & 0.466 & 0.991 & 0.281 & 0.745 \\
Qwen-2.5-7B     & 0.442 & 0.919 & 0.247 & 0.822 \\
Mistral-7B-v0.3 & 0.449 & 0.884 & 0.283 & 0.918 \\
Yi-1.5-9B       & 0.387 & 0.840 & 0.229 & 0.763 \\
\hline
\end{tabular}
\end{table}

\paragraph{Merging methods.} We merge the four adapters with
uniform weights using five methods: Task Arithmetic (TA;
structure-blind averaging), TIES, DARE (both at density $0.2$),
KnOTS\footnote{We run KnOTS's linear inner-merge (shared-subspace
projection then average), not the full SVD-tie alignment of the
original; our KnOTS$\,\approx\,$TA finding is therefore stated for this
variant, though at the zero subspace overlap we measure the theory
predicts no method that exploits overlap can gain, SVD-tie included.},
and task-vector quantization (TVQ;
uniform scalar quantization at $b \in \{1,2,4,8,16,32\}$ bits
per parameter). This is $10$ merge configurations per model,
$40$ in total.

\paragraph{Finding 1: a non-vanishing worst-task excess on every
architecture (and the theoretical floor is zero).}
At the no-quantization rate ($b=32$), worst-task excess is
strictly positive and bootstrap-significant for all four models:
Llama $0.218$ $[0.213, 0.224]$, Mistral $0.146$ $[0.141, 0.151]$,
Qwen $0.110$ $[0.106, 0.114]$, Yi $0.099$ $[0.095, 0.102]$
(Task Arithmetic; the other structure-blind methods agree to
within $0.003$). The bound predicts a floor of the form
$B^2(1 - \deff/(Tr))$. Computing $\deff = \rank(\sum_t P_{V_t})$
directly from the trained adapters, we find $\deff = Tr = 64$ in
\emph{every} layer of \emph{every} model: the four task subspaces
are linearly independent, so the predicted floor is exactly zero.
Real merged adapters thus sit \emph{above} the bound's
floor by $0.10$--$0.22$ nats/token. We read this gap as the
suboptimality of current merging methods relative to the
information-theoretic optimum, rather than as a violation: the
bound is a lower envelope over the worst-case distribution
$P^\star$, and real task vectors are more benign than $P^\star$,
leaving headroom that better algorithms could close.

\paragraph{Finding 2: merging difficulty is governed by training,
not attention design.} The worst-task excesses order the models
\textbf{Llama $>$ Mistral $>$ Qwen $>$ Yi} (hardest to easiest;
all adjacent gaps have non-overlapping $95\%$ CIs). This ordering
cuts across attention design: the GQA model (Qwen) is neither the
easiest nor the hardest, and Yi-1.5-9B -- which shares
Llama-3.1's MHA architecture -- is $2.2\times$ easier to merge
than Llama-3.1 itself. A natural hypothesis is that the ordering
reflects subspace overlap. We test this with two effective-dimension
statistics of the stacked subspace
$M = [V_1 \mid \cdots \mid V_T]$: the hard rank $\deff$ (above,
saturated at $Tr$ for all models) and a soft effective dimension
$\deff^{\mathrm{soft}} = (\sum_i \sigma_i^2)^2 / \sum_i \sigma_i^4$
(participation ratio of the singular values $\sigma_i$ of $M$).
The soft statistic is nearly identical across architectures
($16.4 \pm 0.1$, i.e.\ within $0.7\%$), while worst-task excess
varies by $2.2\times$; the rank correlation between
$\deff^{\mathrm{soft}}/(Tr)$ and worst-task excess is not
significant ($\rho = 0.4$, $p = 0.6$). \emph{Neither
subspace-overlap statistic explains the ordering.} We conclude
that merging-readiness is a property of the pretraining and
instruction-tuning recipe -- plausibly mediated by the curvature
$H_t$ of the loss, which the bound treats as given -- and not of
attention design or subspace geometry alone. Characterizing this
dependence is left to future work.

\paragraph{Finding 3: structure-respecting merging helps only via
the bottleneck task.} TIES is the only method that meaningfully
improves on Task Arithmetic: it reduces worst-task excess by
$30\%$ (Llama), $62\%$ (Mistral), $54\%$ (Yi), and $88\%$ (Qwen)
against the $3$-seed means (\S\ref{subsec:exp-multiseed-bootstrap}),
with non-overlapping CIs in every case. A paired bootstrap of the
per-task differences (the shared $\tau_t$ baseline cancels) shows
the win is concentrated on \emph{the bottleneck task}: on GSM8K,
TIES beats TA by $0.058$--$0.106$ nats on all four models, while
on the easier tasks it is mixed and occasionally significantly
worse. TIES's sign-election and magnitude-pruning resolve
destructive interference where it matters most -- on the
worst-case task -- at a small cost elsewhere. By contrast, KnOTS
and DARE are statistically and practically indistinguishable from
TA (paired differences $\leq 0.002$ and $\leq 0.003$ nats
respectively, split in both directions). DARE's near-identity to TA
is expected rather than a no-op bug: at drop rate $0.2$ with the
standard inverse-keep rescaling, DARE is an \emph{unbiased} estimate
of the TA task vector, so a single draw concentrates near TA on
continuous NLL (its paired differences are nonzero but sign-split) and
can round \emph{exactly} equal on the discrete pass@1 metric of
\S\ref{subsec:exp-downstream-gsm8k}. This is consistent with
the theory: KnOTS's subspace alignment is predicted to help only
when subspaces overlap ($\deff < Tr$), and we measure $\deff = Tr$
everywhere, so there is no overlap for it to exploit. Adopting a
practical-significance threshold of $0.005$ nats, only TIES
separates from the structure-blind cluster
(Figure~\ref{fig:methods}); this rests on our own measurements (see
also~\cite{concurrent2026merging}, a concurrent code-analysis study,
which we cite for context rather than as corroboration).

\paragraph{Finding 4: a universal ``less is more'' rate at
$b=2$.} The TVQ rate sweep is non-monotone on every model: $2$-bit
quantization yields a sharp local minimum in worst-task excess,
$1.7\times$--$6.0\times$ below the adjacent $b=4$ rate (ratios:
Yi $1.65\times$, Llama $2.09\times$, Mistral $2.33\times$, Qwen
$5.97\times$; all with non-overlapping CIs against $b=4$), while
$b=1$ (too coarse) and $b\geq 4$ (too fine) both regress to the
$\approx$TA floor (Figure~\ref{fig:tvq}). On Qwen the
$b=2$ \emph{average}-task excess is negative ($-0.007$): the
$2$-bit merge is better than the mean task adapter. The bound does
not predict this dip -- it upper-bounds quantization-induced
distortion and does not preclude quantization being
\emph{beneficial} -- but the effect is robust across all four
architectures, suggesting that coarse quantization at $b=2$ acts
as a regularizer that suppresses cross-task interference.

We offer one speculative mechanism, which the paper's own
machinery makes natural. Per-tensor uniform quantization at
$b=2$ uses four levels with the decision boundary nearest zero
mapping the smallest-magnitude coordinates to $0$; since merged
task vectors are heavy-tailed with most mass near zero, this
zeroes out roughly the half of coordinates closest to the origin.
That is precisely the trim-then-average operation TIES performs
explicitly -- $b=2$ TVQ is, in effect, an \emph{implicit} TIES,
pruning the low-magnitude coordinates that carry most of the
destructive cross-task interference while preserving the dominant
directions. This predicts (i) the dip should track the sparsity
of the merged delta, and (ii) $b=2$ TVQ and TIES should help on
the same tasks -- both testable, and consistent with TIES being
the only other method that beats Task Arithmetic here
(Finding~3). We flag a full mechanistic account as open, but note
the explanation is not free-floating: it is the same
interference-pruning principle the theory and TIES already rely
on.

\begin{figure}[t]
\centering
\includegraphics[width=0.7\textwidth]{paper_artifacts/figures/headline_methods_ci.png}
\caption{Non-TVQ merging methods, four architectures, $95\%$
bootstrap CIs. Only TIES separates from the Task-Arithmetic /
DARE / KnOTS cluster; KnOTS and DARE are within practical noise
of Task Arithmetic on every model.}
\label{fig:methods}
\end{figure}

\begin{figure}[t]
\centering
\includegraphics[width=0.7\textwidth]{paper_artifacts/figures/headline_tvq_ci.png}
\caption{TVQ rate sweep, four architectures, $95\%$ bootstrap
CIs. A universal local minimum at $b=2$ sits well below the
adjacent rates on every model; $b=1$ (too coarse) and $b\geq 4$
(too fine) regress to the $\approx$Task-Arithmetic floor.}
\label{fig:tvq}
\end{figure}

\paragraph{Scope notes.}
(i) The $2^{-2R/\deff}$ rate-decay slope is operationally invisible
above sub-bit precision in this regime: for $b \geq 1$ the
merging-geometry error dominates the quantization error, leaving at
most one or two TVQ points above the floor on any model. The decay
term is real but not separately resolvable at the rates we measure.
(ii) Translation shows consistently negative excess (merging
\emph{improves} it relative to the translation-only adapter), which
we attribute to under-training of the single-task translation
adapter ($7500$ WMT19 pairs). This does not affect the headline
worst-task metric since GSM8K is the bottleneck on every model, but
the per-task translation row should be read with caution. The full
set of scope items including the
translation under-training caveat is enumerated in
\S\ref{sec:discussion}.

\paragraph{Closing the $0.10$--$0.22$ gap.} The achievability
construction of \S\ref{sec:achv} is what
\S\ref{subsec:exp-rd-encoder} now runs on the real adapters; it
closes the gap to $0.094$ on Llama-3.1-8B-Instruct ($3$-seed mean;
Figure \ref{fig:salvage}), and the regularized variant has the lowest
worst-task excess of every method at its default hyperparameters on
every base we tested.

\paragraph{Seed robustness, now measured
(\S\ref{subsec:exp-multiseed-bootstrap}).} The CIs above quantify
evaluation, not training, uncertainty. The
floor finding ($\deff = Tr$) is a \emph{rank} statement about
$T=4$ rank-$16$ subspaces (a $64$-dimensional span) inside
$\geq 4096$-dimensional layers. We are careful here: \emph{random}
$64$-dimensional subspaces in a $\geq 4096$-dimensional space are in
general position almost surely, but trained LoRA subspaces are
\emph{optimized}, not random, so $\deff = Tr$ is an \emph{empirical}
finding --- we observe it in $100\%$ of layers across four
independently pretrained models --- not a theorem. The large ambient
dimension makes accidental overlap unlikely, which is why we expect it
to be seed-stable, but we do not dress it as inevitable. The model ordering and the $b=2$ dip are
large-margin effects (adjacent-model floor gaps and the
$b{=}2$-vs-$b{=}4$ ratios all have non-overlapping $95\%$ CIs,
\S\ref{subsec:exp-lora}), so a single re-seed cannot overturn them.
We confirm this directly in \S\ref{subsec:exp-multiseed-bootstrap}:
retraining the $T=4$ matrix at \emph{three} independent seeds leaves
all $20$ method-pair gaps that drive our recommendations
\emph{Confident} (sign-stable, $95\%$ CI bounded away from zero).

\paragraph{Two adapter cohorts; reading the numbers across
\S\ref{sec:experiments}.} Two cohorts recur below and are \emph{not}
cross-comparable. The \emph{$T=4$ matrix} (GSM8K, Alpaca, Magicoder,
FLORES translation) underlies
\S\ref{subsec:exp-rd-encoder}--\S\ref{subsec:exp-downstream-gsm8k}
and \S\ref{subsec:exp-e10}; the \emph{$T$-scaling pool} (those four
plus CodeAlpaca, Dolly, XSum) underlies \S\ref{subsec:exp-T-scaling},
where merges draw $T \in \{2,4,7\}$ subsets. A base's worst-task
excess differs between cohorts because the task set differs: Llama
Task Arithmetic is $\approx 0.22$ on the $T=4$ matrix but $0.153$ on
a $T=4$ pool subset (\S\ref{subsec:exp-T-scaling}), and the two must
not be cross-read. Within the $T=4$ matrix, single-seed point
estimates (quoted in \S\ref{subsec:exp-rd-encoder},
\S\ref{subsec:exp-downstream-gsm8k}, \S\ref{subsec:exp-e10}) and the
$3$-seed means (\S\ref{subsec:exp-multiseed-bootstrap}) agree to
within a $\pm 0.01$ seed band on the point estimates (widest for
TVQ$_2$, whose seed-1 cell sits $\sim 0.009$ below its $3$-seed mean;
Llama TA spans $0.213$--$0.220$, TIES $0.147$--$0.161$), while the
\emph{method-pair gaps} that drive the recommendations are tighter
($\pm 0.004$, \S\ref{subsec:exp-multiseed-bootstrap}). The multi-seed
bootstrap of \S\ref{subsec:exp-multiseed-bootstrap} is the canonical
reference for the matrix.
